\documentclass[11pt]{article}

% Shared preamble for Thurstone papers.
% Imported via % Shared preamble for Thurstone papers.
% Imported via % Shared preamble for Thurstone papers.
% Imported via \input{../shared/preamble.tex} from each paper.tex.
% Keep packages and macros common to all papers here.

\usepackage[utf8]{inputenc}
\usepackage[T1]{fontenc}
\usepackage{lmodern}
\usepackage{amsmath,amssymb,amsthm}
\usepackage{graphicx}
\graphicspath{{figures/}}
\usepackage{booktabs}
\usepackage{microtype}
\usepackage[hidelinks]{hyperref}

% Bibliography. We use biblatex with the bibtex backend because Tectonic runs
% bibtex natively (no external biber, which avoids biber/biblatex version
% mismatches). bibtex runs through Tectonic's own I/O, so the shared
% bibliography at papers/refs.bib resolves directly via the parent path below.
\usepackage[backend=bibtex,style=authoryear,sorting=nyt]{biblatex}
\addbibresource{../refs.bib}

% --- Shared macros ---
\newcommand{\thurstone}{Thurstone}
 from each paper.tex.
% Keep packages and macros common to all papers here.

\usepackage[utf8]{inputenc}
\usepackage[T1]{fontenc}
\usepackage{lmodern}
\usepackage{amsmath,amssymb,amsthm}
\usepackage{graphicx}
\graphicspath{{figures/}}
\usepackage{booktabs}
\usepackage{microtype}
\usepackage[hidelinks]{hyperref}

% Bibliography. We use biblatex with the bibtex backend because Tectonic runs
% bibtex natively (no external biber, which avoids biber/biblatex version
% mismatches). bibtex runs through Tectonic's own I/O, so the shared
% bibliography at papers/refs.bib resolves directly via the parent path below.
\usepackage[backend=bibtex,style=authoryear,sorting=nyt]{biblatex}
\addbibresource{../refs.bib}

% --- Shared macros ---
\newcommand{\thurstone}{Thurstone}
 from each paper.tex.
% Keep packages and macros common to all papers here.

\usepackage[utf8]{inputenc}
\usepackage[T1]{fontenc}
\usepackage{lmodern}
\usepackage{amsmath,amssymb,amsthm}
\usepackage{graphicx}
\graphicspath{{figures/}}
\usepackage{booktabs}
\usepackage{microtype}
\usepackage[hidelinks]{hyperref}

% Bibliography. We use biblatex with the bibtex backend because Tectonic runs
% bibtex natively (no external biber, which avoids biber/biblatex version
% mismatches). bibtex runs through Tectonic's own I/O, so the shared
% bibliography at papers/refs.bib resolves directly via the parent path below.
\usepackage[backend=bibtex,style=authoryear,sorting=nyt]{biblatex}
\addbibresource{../refs.bib}

% --- Shared macros ---
\newcommand{\thurstone}{Thurstone}


\usepackage{amsmath}
\usepackage{algorithm}
\usepackage{algpseudocode}
\usepackage{cleveref}

\title{Thurstone Portfolios:\\
       Long-Only Allocation by Inverting Winning Probabilities}
\author{Peter Cotton}
\date{\today}

\newcommand{\E}{\mathbb{E}}
\newcommand{\Prob}{\mathbb{P}}
\newcommand{\R}{\mathbb{R}}
\DeclareMathOperator*{\argminop}{arg\,min}

\begin{document}
\maketitle

\begin{abstract}
We develop \emph{Thurstone portfolios}, a long-only asset-allocation scheme in
which portfolio weights are the winning probabilities of a Thurstonian race
among assets. Starting from a variance-only (diagonal) allocation, we invert the
weights into latent \emph{abilities} using the calibration map of
\textcite{cotton2021}, re-run the contest with correlation reintroduced, and
read off each asset's probability of being the ``winner'' as its new weight. The
construction is long-only and budget-feasible by design, requires no covariance
inversion, and degrades gracefully when the covariance matrix is ill-conditioned
or singular. A single scaling parameter $\phi$ interpolates between ignoring
correlation and using it in full, giving a one-dimensional dial that plays the
role of---but is more stable than---shrinkage in mean--variance optimization. We
position the method against the Capital Asset Pricing Model (CAPM), whose
implied market-capitalization weighting we argue is a \emph{Lucian} (proportional
choice) special case that a Thurstonian model need not, and empirically does not,
reproduce. We outline a backtesting protocol and report on two empirical studies
(REITs and Russell-index equities tilted by language-model ``AI-benefit''
scores) used to evaluate the approach.
\end{abstract}

\section{Introduction}
\label{sec:intro}

Modern passive investing rests on an equilibrium argument: under the CAPM
\parencite{sharpe1964}, the market-capitalization-weighted portfolio is mean--variance
efficient, so a rational investor should simply hold the index. There is a less
remarked structural assumption underneath. Capitalization weighting treats the
relative ``attractiveness'' of two assets as a fixed ratio that is preserved
when other assets are added to or removed from the universe. This is precisely
Luce's choice axiom \parencite{luce1959}, the independence-of-irrelevant-alternatives
property that underlies the Bradley--Terry \parencite{bradleyterry1952} and
Plackett--Luce \parencite{plackett1975} models of choice and ranking. We will call a
weighting scheme with this property \emph{Lucian}.

The Lucian assumption is convenient but neither logically necessary nor
empirically robust. Investors do not observe the true covariance of returns;
they estimate it, and good practice is to \emph{shrink} those estimates toward a
structured target \parencite{ledoitwolf2004}. Shrinkage breaks the
wisdom-of-crowds aggregation that would make the capitalization-weighted
portfolio efficient, and tracking the index then carries a cost in
risk-adjusted return unless the index is itself optimal \parencite{roll1992}.
This invites a question: if the market need not reveal itself as exactly Lucian,
what is a natural non-Lucian alternative?

Thurstone's law of comparative judgment \parencite{thurstone1927} supplies one.
In a Thurstonian model each asset $i$ carries a latent \emph{ability} $a_i$ and a
noisy \emph{performance} $X_i = a_i + \varepsilon_i$; the asset with the extreme
performance ``wins.'' Unlike the Lucian model, the probability that $i$ beats $j$
depends on the rest of the field and---crucially for finance---on the
\emph{correlation} of the performances. The contribution of this paper is to use
the Thurstonian winning probabilities themselves as portfolio weights, and to
show that the inverse problem (recovering abilities from a target allocation, and
re-deriving weights once correlation is injected) admits a fast, stable solution.

\paragraph{Contributions.}
\begin{enumerate}
  \item We define the \emph{ability tilt}: a map from a baseline diagonal
        portfolio to a correlation-aware long-only portfolio, via inversion to
        Thurstone abilities and a correlated race (\Cref{sec:method}).
  \item We give the construction both as a Monte-Carlo / quasi-Monte-Carlo
        sampler (as implemented in the \texttt{winningport} package) and as a
        deterministic lattice computation using the \texttt{thurstone} package's
        order-statistics machinery (\Cref{sec:computation}).
  \item We articulate the \emph{Indispensable Markets Hypothesis} as a
        non-Lucian alternative to the CAPM and locate capitalization weighting
        as its degenerate, zero-correlation limit (\Cref{sec:capm}).
  \item We specify a backtesting protocol and summarize two empirical studies
        used to evaluate the method (\Cref{sec:experiments}).
\end{enumerate}

\section{The Thurstonian Model of a Field}
\label{sec:model}

Consider a universe of $n$ assets. Each asset has a latent ability $a_i\in\R$ and
a performance
\begin{equation}
  X_i = a_i + \varepsilon_i, \qquad (\varepsilon_1,\dots,\varepsilon_n) \sim
  \mathcal{N}(0, \Sigma_\varepsilon),
  \label{eq:performance}
\end{equation}
where we adopt the convention that the \emph{minimum} performance wins, so that a
smaller ability denotes a stronger competitor. The \emph{state price} (winning
probability) of asset $i$ is
\begin{equation}
  p_i \;=\; \Prob\!\left( X_i = \min_{k} X_k \right).
  \label{eq:stateprice}
\end{equation}
When the performances are independent ($\Sigma_\varepsilon$ diagonal) the
$p_i$ are determined by the abilities alone through the distribution of the first
order statistic, and $p$ is a smooth, monotone function of $a$
\parencite{cotton2021}. Two maps are then of interest:
\begin{align}
  \text{forward:}\quad & a \;\longmapsto\; p, \label{eq:forward}\\
  \text{inverse:}\quad & p \;\longmapsto\; a. \label{eq:inverse}
\end{align}
The inverse \eqref{eq:inverse}---the \emph{Fast Ability Transform}---recovers
abilities (up to an additive constant) from a target set of winning
probabilities, and is the engine of the construction below.

\section{Thurstone Portfolios via the Ability Tilt}
\label{sec:method}

We treat portfolio \emph{weights} as winning probabilities and \emph{assets} as
competitors. The procedure ``polishes'' a baseline allocation by reintroducing
correlation through a Thurstonian race.

\begin{algorithm}[h]
\caption{Ability tilt (Thurstone portfolio)}
\label{alg:tilt}
\begin{algorithmic}[1]
\Require Covariance estimate $\Sigma\in\R^{n\times n}$; correlation scale
         $\phi\in[0,1]$.
\State $w \gets \textsc{DiagonalPortfolio}(\Sigma)$
       \Comment{variance-only long-only weights}
\State $a \gets \textsc{StatePriceImpliedAbility}(w)$
       \Comment{inverse map \eqref{eq:inverse}}
\State $C \gets \textsc{Corr}(\Sigma)$; \quad
       $C_\phi \gets \textsc{ScaleOffDiagonal}(C, \phi)$
\State $C_\phi \gets \textsc{NearestValidCorrelation}(C_\phi)$
\State Draw correlated performances $X^{(1)},\dots,X^{(M)} \sim
       \mathcal{N}(a, C_\phi)$
\State $w'_i \gets \dfrac{1}{M}\sum_{m=1}^{M}
       \mathbf{1}\!\left[\, i = \argminop_k X^{(m)}_k \,\right]$
       \Comment{win frequency}
\State \Return $w'$
\end{algorithmic}
\end{algorithm}

Several properties follow directly from \Cref{alg:tilt}:
\begin{itemize}
  \item \textbf{Long-only and budget-feasible.} Each $w'_i \ge 0$ and
        $\sum_i w'_i = 1$ because the weights are empirical winning frequencies.
        No constraints or projections are needed.
  \item \textbf{No covariance inversion.} The method never inverts $\Sigma$. It
        needs only a valid correlation matrix for sampling, repaired if
        necessary by projection to the nearest correlation matrix. This is what
        lets it tolerate ill-conditioned or singular $\Sigma$.
  \item \textbf{One interpretable dial.} At $\phi = 0$ the off-diagonals vanish
        and the race recovers (a Monte-Carlo estimate of) the diagonal portfolio;
        at $\phi = 1$ the full estimated correlation is used. Intermediate $\phi$
        plays the role of shrinkage, but acts on the \emph{correlation injected
        into the race} rather than on the covariance estimate itself.
  \item \textbf{Diversification pressure.} An asset that is highly correlated
        with strong competitors wins the minimum less often once correlation is
        switched on, so the tilt down-weights redundant exposures and up-weights
        diversifiers---without an explicit risk penalty.
\end{itemize}

\section{Computation}
\label{sec:computation}

\paragraph{Sampling implementation.}
The \texttt{winningport} package realizes \Cref{alg:tilt} with quasi-Monte-Carlo
sampling: performances are drawn with a randomized Sobol' sequence
(\texttt{MultivariateNormalQMC}) and weights are the argmin frequencies. The
sample budget $M$ is rounded to a power of two for the Sobol' construction, and
$\phi$ scales the off-diagonal correlation entries before a
nearest-correlation repair.

\paragraph{Deterministic lattice implementation.}
The \texttt{thurstone} package computes winning probabilities exactly on a
discrete lattice rather than by sampling, building the distribution of the first
order statistic by convolving competitor densities and tracking conditional
multiplicity for ties \parencite{cotton2021}. In the independent case this gives
the forward map \eqref{eq:forward} and its inverse \eqref{eq:inverse} without
Monte-Carlo error and in time linear in $n$. Extending the deterministic path to
the correlated race (e.g.\ via a Gaussian-copula common-factor reduction) is the
natural bridge between the legacy sampling code and the current package, and we
treat it as the reference implementation for reproducibility.

\section{The Indispensable Markets Hypothesis}
\label{sec:capm}

The CAPM implies capitalization weighting, which is Lucian: relative weights are
fixed ratios independent of the rest of the universe and of correlation. The
ability tilt at $\phi = 0$ inherits this independence and is, in that limit, a
Lucian allocator. Turning $\phi$ up makes the allocator non-Lucian: weights now
respond to the correlation structure of the field.

We state the \emph{Indispensable Markets Hypothesis} informally as the claim that
markets reveal, through prices, a Thurstonian ability structure that is
\emph{not} recoverable from capitalization alone, and that a portfolio built from
inferred abilities and the correlation-aware race can earn a higher long-run
risk-adjusted return than the capitalization-weighted index. The hypothesis is
falsifiable: it predicts that the advantage of the tilt over capitalization
weighting grows with the heterogeneity of pairwise correlations and vanishes as
the correlation structure flattens. \Cref{sec:experiments} describes the tests.

\section{Empirical Protocol and Studies}
\label{sec:experiments}

\paragraph{Protocol.}
For each rebalancing date we estimate $\Sigma$ from a trailing window, form the
diagonal portfolio, apply \Cref{alg:tilt} at a grid of $\phi$ values and sample
budgets $M$, and hold the resulting weights until the next rebalance. We report
annualized return, an adjusted Sharpe ratio, and maximum drawdown, benchmarked
against equal-weight and capitalization-weight portfolios.

\paragraph{Study 1: REITs.}
A panel of real-estate investment trusts (daily returns and market caps) is used
to compare the ability tilt against equal- and capitalization-weighted
baselines across sample budgets $M\in\{10^3,10^5\}$ and correlation scales
$\phi\in\{0,\tfrac12,1\}$. \emph{Results to be inserted.}

\paragraph{Study 2: Russell indices with AI-benefit tilts.}
Russell 1000/2000/3000 constituents are scored by language models for expected
benefit from AI adoption; the scores define the field that is then allocated by
the ability tilt and compared with capitalization weighting and with simple
top-quantile score tilts. \emph{Results to be inserted.}

\begin{table}[h]
  \centering
  \caption{Backtest summary (placeholder).}
  \label{tab:results}
  \begin{tabular}{lccc}
    \toprule
    Portfolio & Ann.\ return & Adj.\ Sharpe & Max drawdown \\
    \midrule
    Equal weight        & --- & --- & --- \\
    Cap weight          & --- & --- & --- \\
    Diagonal ($\phi{=}0$) & --- & --- & --- \\
    Ability tilt ($\phi{=}\tfrac12$) & --- & --- & --- \\
    Ability tilt ($\phi{=}1$) & --- & --- & --- \\
    \bottomrule
  \end{tabular}
\end{table}

\section{Discussion}
\label{sec:discussion}

The ability tilt is a small idea with a few useful properties: it is long-only
and feasible without optimization, it avoids covariance inversion, and it exposes
a single correlation dial in place of the many knobs of constrained
mean--variance optimization. Open questions include the choice and stability of
$\phi$ over time, the deterministic correlated-race computation, the treatment of
estimation error in abilities (a Thurstonian analogue of covariance shrinkage),
and turnover once the method is rebalanced on real data.

\section{Conclusion}
\label{sec:conclusion}

Thurstone portfolios weight assets by their probability of winning a correlated
race among latent abilities inferred from a baseline allocation. The construction
is a natural non-Lucian alternative to capitalization weighting and is cheap and
robust to compute. Whether it pays for itself out of sample is an empirical
question, and the protocol and studies above are designed to answer it.

\printbibliography

\end{document}
